% Source: https://artofproblemsolving.com/community/c3h1600483_sine_cosine_tangent
% Stable identity: aops-topic-1600483-q01
Each value is a ratio of sides of a right triangle, so each angle is found with the inverse function (a calculator's $\sin^{-1}$, $\cos^{-1}$, $\tan^{-1}$ in degree mode, or a table of values):

\begin{asy}
usepackage("amsmath");
size(11cm);
defaultpen(fontsize(10pt));
real d2r=pi/180;real k=3;
real[] th={15.5,29.5,24.5};
pair[] o={(0,0),(0,-2.3),(0,-4.6)};
for(int i=0;i<3;++i){
  pair A=o[i];pair B,C;
  if(i<2){B=A+k*(cos(th[i]*d2r),0);C=A+k*(cos(th[i]*d2r),sin(th[i]*d2r));}
  else {B=A+(k,0);C=A+(k,k*0.4552);}
  draw(A--B--C--cycle,linewidth(0.9));
  draw(B+(-0.15,0)--B+(-0.15,0.15)--B+(0,0.15));
  draw(arc(A,0.9,0,th[i]),blue);
  label("$\theta_"+string(i+1)+"$",A+1.25*dir(th[i]/2),blue);
}
label("$1$",o[0]+1.7*dir(15.5),dir(105.5));label("$0.2672$",o[0]+(k*cos(15.5*d2r),k*sin(15.5*d2r)/2),E);
label("$1$",o[1]+1.7*dir(29.5),dir(119.5));label("$0.8704$",o[1]+(k*cos(29.5*d2r)/2,0),S);
label("$1$",o[2]+(k/2,0),S);label("$0.4552$",o[2]+(k,k*0.4552/2),E);
label("\mbox{}",(min(currentpicture,user=true)+max(currentpicture,user=true))/2,black); // keeps the text after the figure black when LaTeX typesets the labels inline
\end{asy}

\begin{enumerate}
\item[(a)] $\theta_1=\sin^{-1}(0.2672)\approx15.498^\circ$, so $\theta_1\approx\boxed{15.5^\circ}$. Check: $\sin15.5^\circ\approx0.26724$.
\item[(b)] $\theta_2=\cos^{-1}(0.8704)\approx29.495^\circ$, so $\theta_2\approx\boxed{29.5^\circ}$. Check: $\cos29.5^\circ\approx0.87036$.
\item[(c)] $\theta_3=\tan^{-1}(0.4552)\approx24.475^\circ$, so $\theta_3\approx\boxed{24.5^\circ}$. Check: $\tan24.5^\circ\approx0.45573$, and $\tan24.4^\circ\approx0.45362$, so $0.4552$ lies between them, closer to $24.5^\circ$.
\end{enumerate}
In a table, look up the value in the sine (cosine, tangent) column and read off the angle beside it, interpolating between neighbouring entries when the value is not listed exactly.
